% answers_short.tex —— 1991 年·全国卷·理科 · 简短答案速查
% 由 tools/make_short_answers.py 自 answers.tex 的 \gksolhead/\gkanswer 提取，
% 并与 1991/1991_数学真题_全国理_答案要点.md 逐题比对后生成，勿手改（改源 MD 后重跑）。

\begin{gknotice}
\textbf{使用说明}：本节只给最终答案便于快速核对，完整推导见解析版。答案为编写者独立推导，\textbf{非官方标准答案}。
\end{gknotice}

\gkpart{选择题}
\begin{center}\begin{tabular}{cccccccccccccccc}
\toprule
题号 & 1 & 2 & 3 & 4 & 5 & 6 & 7 & 8 & 9 & 10 & 11 & 12 & 13 & 14 & 15 \\
\midrule
答案 & A & D & B & B & A & D & A & D & C & C & A & C & B & C & D \\
\bottomrule
\end{tabular}\end{center}

\gkpart{填空题}
\begin{itemize}[leftmargin=2.6em]
  \item[\textbf{16.}] $\dfrac\pi4$
  \item[\textbf{17.}] $(-2,1)$
  \item[\textbf{18.}] $\dfrac{14}3$
  \item[\textbf{19.}] $1+\dfrac{\sqrt{10}}5$
  \item[\textbf{20.}] $3\pi a^2$
  \item[\textbf{21.}] 最小值为 $2-\sqrt2$；$x=-\dfrac{3\pi}8+k\pi$，$k\in\mathbb Z$。
  \item[\textbf{22.}] 模为 $\sqrt2$，辐角的主值为 $-\dfrac\pi4$。
\end{itemize}

\gkpart{解答题}
\begin{itemize}[leftmargin=2.6em]
  \item[\textbf{23.}] $\dfrac2{\sqrt{11}}$。
  \item[\textbf{24.}] $f(x)=-x^3+1$ 在 $\mathbb R$ 上严格递减。
  \item[\textbf{25.}] 原式要求 $n\ge1$，定义域为 $x>\sqrt a$。若 $n$ 为奇数，则 $\sqrt a<x<\dfrac{1+\sqrt{1+4a}}2$；若 $n$ 为偶数，则 $x>\dfrac{1+\sqrt{1+4a}}2$。
  \item[\textbf{26.}] $x^2-\dfrac{y^2}3=1$。
\end{itemize}
